\documentclass[11pt,letterpaper]{article}

\usepackage[spanish,es-noshorthands]{babel}
\usepackage{fontspec}
\setmainfont{Source Serif Pro}
\usepackage{microtype}
\usepackage{geometry}
\usepackage{booktabs}
\usepackage{tabularx}
\usepackage{enumitem}
\usepackage{xcolor}
\usepackage{graphicx}
\usepackage{csquotes}
\usepackage{amsmath}
\usepackage{siunitx}
\usepackage{tikz}
\usetikzlibrary{arrows.meta,positioning,shapes.geometric,fit,backgrounds,patterns}
\usepackage{xurl}
\usepackage[hidelinks]{hyperref}
\usepackage[backend=biber,style=apa,sorting=nyt]{biblatex}
\DeclareLanguageMapping{spanish}{spanish-apa}

\geometry{margin=2.3cm}
% Los rótulos de las figuras usan nodos de ancho fijo con texto alineado a la
% derecha, y la bibliografía compone en modo sloppy: ambos generan avisos de
% línea holgada que son cosméticos y no indican un problema de composición.
% Se silencian para que `make check` siga siendo sensible a los desbordes reales.
\hbadness=10000
\setlength{\emergencystretch}{1em}
\setlength{\parindent}{0pt}
\setlength{\parskip}{0.55em}
\setlist{nosep}
\addbibresource{referencias.bib}
\AtBeginBibliography{\sloppy}

\sisetup{
  output-decimal-marker = {,},
  group-separator = {.},
  group-minimum-digits = 4,
  detect-all
}

% Colores de las cuatro generaciones tecnológicas
\definecolor{cImpresa}{RGB}{31,78,121}
\definecolor{cEncarta}{RGB}{191,110,30}
\definecolor{cWiki}{RGB}{40,124,70}
\definecolor{cLLM}{RGB}{116,66,150}

% ---------------------------------------------------------------------------
% Perfil de figuras de mérito: eje logarítmico común de 10^-3 a 10^9.
% Posición horizontal: x = 0,92 * (log10(valor) + 3)
% ---------------------------------------------------------------------------
\newcommand{\fomgrid}[1]{%
  \foreach \e in {-3,-2,...,9} {%
    \pgfmathsetmacro{\xg}{0.92*(\e+3)}%
    \draw[gray!16] (\xg,-0.45) -- (\xg,#1);%
  }%
  \foreach \e/\lab in {-3/{$10^{-3}$}, -1/{$10^{-1}$}, 1/{$10^{1}$},
                       3/{$10^{3}$}, 5/{$10^{5}$}, 7/{$10^{7}$}, 9/{$10^{9}$}} {%
    \pgfmathsetmacro{\xg}{0.92*(\e+3)}%
    \draw[gray!45] (\xg,-0.45) -- (\xg,#1);%
    \node[below, font=\tiny, gray!55!black] at (\xg,-0.45) {\lab};%
  }%
}

\newcommand{\fombar}[5]{%
  \node[left, text width=3.5cm, align=right, font=\scriptsize] at (-0.12,#1) {#4};%
  \fill[#3] (0,#1-0.17) rectangle (#2,#1+0.17);%
  \node[right, font=\scriptsize] at (#2+0.10,#1) {#5};%
}

% Rótulos de fila comunes a las cuatro fichas de FOM
\newcommand{\rCob}{Cobertura\\\tiny (artículos)}
\newcommand{\rPro}{Profundidad\\\tiny (palabras por artículo)}
\newcommand{\rLat}{Latencia de actualización\\\tiny (segundos)}
\newcommand{\rAcc}{Latencia de acceso\\\tiny (segundos)}
\newcommand{\rCos}[1]{Coste marginal\\por #1\\\tiny (dólares)}
\newcommand{\notafom}{En las tres últimas filas, un valor menor indica mejor
  desempeño.}

\title{\vspace{-1.2cm}\textbf{Tecnologías invasoras: de la enciclopedia impresa
a los modelos de lenguaje}\\
\large Cuatro curvas en S sobre un mismo eje de mérito, 1768--2026}
\author{Carlos Luis Rojas Aragonés\\
\small Gestión del Desarrollo Tecnológico}
\date{\small 14 de agosto de 2026}

\begin{document}

\maketitle
\vspace{-0.6em}

\section{Introducción}

El caso que analizo es la sustitución de \textbf{la enciclopedia impresa}, con la
\emph{Encyclopædia Britannica} como testigo, publicada de forma continuada desde
1768 en Edimburgo hasta que el 13 de marzo de 2012 la compañía anunció que
dejaba de imprimirla, después de 244 años \parencite{britannica2012,npr2012}.

El relato popular dice que la mató Wikipedia. Los datos dicen algo más
interesante y más útil para la gestión tecnológica: no hubo un relevo entre dos
tecnologías, sino una \textbf{cadena de cuatro}, y cada eslabón ganó en una
figura de mérito distinta de la que el titular sabía optimizar.

\begin{enumerate}
  \item \textbf{La enciclopedia impresa} (1768--2012): el titular. Insuperable en
    profundidad y autoridad, prisionera del papel en latencia y coste.
  \item \textbf{Microsoft Encarta} (1993--2009): el primer invasor, sobre soporte
    CD-ROM. Peor en contenido, mejor en coste y en acceso.
  \item \textbf{Wikipedia} (2001 hasta hoy): el segundo invasor, sobre la web y
    el modelo wiki. Ganó primero en cobertura y después en casi todo lo demás.
  \item \textbf{Los modelos de lenguaje} (2022 hasta hoy): el tercer invasor.
    Redefine la pregunta: ya no entrega el artículo, entrega la respuesta.
\end{enumerate}

Hay un dato que ordena todo el análisis y que conviene adelantar: \textbf{Encarta
cerró en octubre de 2009, casi tres años antes de que la Britannica dejara de
imprimirse} \parencite{techcrunch2009}. El primer invasor murió antes que su
víctima, eliminado por el segundo invasor. Un caso de sustitución tecnológica
no es, por tanto, un duelo, sino una secuencia en la que los invasores también
son invadidos, y cada vez más rápido: la impresa resistió 244 años, Encarta 16 y
Wikipedia lleva 25 con su tráfico humano ya en descenso.

\section{Marco de análisis}

Una \textbf{figura de mérito} (FOM) es una medida cuantitativa del desempeño de
una tecnología en la función que presta, independiente de la implementación
concreta que la realiza \parencite{deweck2022}. Cada tecnología recorre en su
FOM una \textbf{curva en S}: arranque lento, mejora rápida y saturación al
acercarse a un límite físico o económico \parencite{foster1986,utterback1994}.

Una \textbf{tecnología invasora} es la que aparece sobre una curva en S distinta
y termina cruzando la del titular. Tres propiedades del proceso, que este caso
ilustra con una nitidez poco común, conviene tener presentes:

\begin{enumerate}
  \item \textbf{El invasor casi nunca gana en el FOM del titular.} Entra por
    debajo en el atributo que el titular considera decisivo y muy por encima en
    otros que el titular no mide \parencite{christensen1997}.
  \item \textbf{El invasor redefine el FOM.} Cuando la comparación deja de ser
    posible en las mismas unidades, la sustitución ya está decidida.
  \item \textbf{El cruce de las curvas precede a la muerte comercial en años.}
    El titular sigue vendiendo mucho después de haber perdido la superioridad
    técnica, porque su activo real no es la tecnología, sino el canal.
\end{enumerate}

\section{Panorama: cuatro curvas sobre un mismo eje}

La Figura~\ref{fig:curvas} representa el proceso completo con un eje de tiempo y
un eje de FOM. Como FOM principal uso la \textbf{cobertura}, es decir el número
de artículos enciclopédicos que el lector puede consultar, porque es el único
atributo medible en las mismas unidades a lo largo de dos siglos y medio.

\begin{figure}[ht]
  \centering
  \begin{tikzpicture}
    % --- rejilla y eje vertical ---
    \foreach \yv/\lab in {0/{$10^{3}$}, 1.3/{$10^{4}$}, 2.6/{$10^{5}$},
                          3.9/{$10^{6}$}, 5.2/{$10^{7}$}} {
      \draw[gray!20] (-0.1,\yv) -- (13.55,\yv);
      \node[left, font=\scriptsize, gray!55!black] at (-0.15,\yv) {\lab};
    }
    \draw[-{Stealth[length=2mm]}, gray!60!black] (-0.1,-0.05) -- (-0.1,5.6);
    \node[rotate=90, font=\scriptsize, gray!55!black, align=center]
      at (-1.15,2.6) {Artículos accesibles (escala log.)};
    % --- eje de tiempo con corte ---
    \draw[gray!60!black] (-0.1,0) -- (3.60,0);
    \draw[-{Stealth[length=2mm]}, gray!60!black] (3.95,0) -- (13.55,0);
    \draw[gray!60!black] (3.68,-0.12) -- (3.80,0.12);
    \draw[gray!60!black] (3.76,-0.12) -- (3.88,0.12);
    \foreach \xv/\lab in {0.295/1889, 1.015/1911, 1.833/1936, 3.076/1974} {
      \draw[gray!60!black] (\xv,0) -- (\xv,-0.13);
      \node[below, font=\tiny, gray!55!black] at (\xv,-0.13) {\lab};
    }
    \foreach \xv/\lab in {5.20/1995, 6.45/2000, 7.70/2005, 8.95/2010,
                          10.20/2015, 11.45/2020, 12.70/2025} {
      \draw[gray!60!black] (\xv,0) -- (\xv,-0.13);
      \node[below, font=\tiny, gray!55!black] at (\xv,-0.13) {\lab};
    }
    \node[font=\tiny, gray!55!black, align=center] at (1.9,-0.62)
      {escala comprimida};
    % --- Britannica impresa ---
    \draw[very thick, cImpresa] (0.295,1.600) -- (1.015,2.083) -- (3.436,2.357)
      -- (3.60,2.357);
    \draw[very thick, cImpresa] (3.95,2.357) -- (9.50,2.357);
    \foreach \x/\y in {0.295/1.600, 1.015/2.083, 3.436/2.357} {
      \fill[cImpresa] (\x,\y) circle (1.7pt);
    }
    \draw[thick, cImpresa] (9.38,2.24) -- (9.62,2.48);
    \draw[thick, cImpresa] (9.62,2.24) -- (9.38,2.48);
    \node[font=\tiny, cImpresa, above right=-0.6mm and -1mm] at (1.015,2.083)
      {\num{40000} (1911)};
    \node[font=\tiny, cImpresa, right] at (4.05,2.50) {\num{65000} artículos};
    % --- Encarta ---
    \draw[thick, cEncarta, dashed] (4.70,1.75)
      .. controls (6.0,2.15) and (7.0,2.30) .. (8.45,2.330);
    \fill[cEncarta] (8.45,2.330) circle (1.7pt);
    \draw[cEncarta!70, thin] (8.45,2.28) -- (8.62,1.92);
    \node[font=\tiny, cEncarta, right] at (8.58,1.83)
      {\num{62000} (2008)};
    \draw[thick, cEncarta] (8.79,2.21) -- (9.03,2.45);
    \draw[thick, cEncarta] (9.03,2.21) -- (8.79,2.45);
    % --- Wikipedia ---
    \draw[very thick, cWiki] (6.71,0) -- (6.89,1.3) -- (7.215,2.6)
      -- (7.7525,3.509) -- (7.9925,3.9) -- (8.3725,4.291) -- (8.855,4.520)
      -- (9.5825,4.683) -- (10.4075,4.809) -- (11.465,4.912)
      -- (12.8025,4.999) -- (13.105,5.017);
    \foreach \x/\y in {6.71/0, 6.89/1.3, 7.215/2.6, 7.7525/3.509, 7.9925/3.9,
                       8.3725/4.291, 8.855/4.520, 9.5825/4.683, 10.4075/4.809,
                       11.465/4.912, 12.8025/4.999} {
      \fill[cWiki] (\x,\y) circle (1.7pt);
    }
    \node[font=\tiny, cWiki, right=0.5mm] at (13.105,5.017)
      {\num{7224454}};
    \node[font=\tiny, cWiki, above left=-0.8mm and -1.2mm] at (6.71,0)
      {\num{1000}};
    % --- cruce ---
    \draw[red!65!black, thick] (7.154,2.357) circle (3.2pt);
    \draw[red!65!black, thin] (7.02,2.50) -- (6.75,2.82);
    \node[font=\scriptsize, red!65!black, align=center, above left=-0.5mm and -1mm]
      at (6.75,2.82) {cruce del FOM\\[-0.3em] \tiny (2002)};
    % --- LLM ---
    \draw[cLLM, dotted, thick] (12.18,0) -- (12.18,4.7);
    \node[font=\tiny, cLLM, rotate=90, above right=0mm and 0.4mm] at (12.18,2.9)
      {ChatGPT, nov. 2022};
    \fill[cLLM] (12.905,3.831) circle (1.7pt);
    \node[font=\tiny, cLLM, below right=-0.2mm and 0.2mm] at (12.905,3.831)
      {\num{885279}};
    \draw[-{Stealth[length=2mm]}, thick, cLLM, dashed]
      (12.98,3.95) -- (13.30,4.83);
    \node[font=\tiny, cLLM, left] at (11.85,4.40) {cobertura no acotada};
    % --- eventos ---
    \foreach \xv/\lab in {%
      4.70/{1993: Encarta},
      6.20/{1999: britannica.com gratuito},
      6.71/{2001: Wikipedia},
      7.7525/{2005: estudio de \emph{Nature}},
      8.9075/{2009: cierre de Encarta},
      9.50/{2012: fin de la impresión},
      12.18/{2022: ChatGPT},
      12.905/{2025: Grokipedia}} {
      \draw[gray!55] (\xv,0) -- (\xv,-0.3);
      \node[rotate=60, anchor=east, font=\tiny, gray!35!black] at (\xv,-0.34)
        {\lab};
    }
    % --- leyenda ---
    \begin{scope}[shift={(3.20,5.35)}]
      \draw[very thick, cImpresa] (0,0) -- (0.45,0);
      \node[right, font=\tiny] at (0.5,0) {impresa};
      \draw[thick, cEncarta, dashed] (1.95,0) -- (2.40,0);
      \node[right, font=\tiny] at (2.45,0) {Encarta};
      \draw[very thick, cWiki] (3.85,0) -- (4.30,0);
      \node[right, font=\tiny] at (4.35,0) {Wikipedia};
      \draw[thick, cLLM, dashed] (5.95,0) -- (6.40,0);
      \node[right, font=\tiny] at (6.45,0) {modelos de lenguaje};
    \end{scope}
  \end{tikzpicture}
  \caption{Curvas en S de las cuatro generaciones sobre el FOM de cobertura.
  Puntos verificados: \num{17000} artículos en la novena edición de la
  Britannica, \num{40000} en la undécima (1910--1911) y \num{65000} en la
  Micropædia de la decimoquinta; \num{62000} en Encarta Premium 2008
  \parencite{encartawiki}; los once hitos de la Wikipedia en inglés desde
  \num{1000} artículos en febrero de 2001 hasta \num{7000000} el 28 de mayo de
  2025, y \num{7224454} el 13 de agosto de 2026 \parencite{wikistats2026}; y
  \num{885279} entradas de Grokipedia el 27 de octubre de 2025
  \parencite{yasseri2025}. Se dibujan de forma esquemática, por falta de una
  serie anual pública, el tramo de la Britannica anterior a 1889 (no
  representado) y el de Encarta entre 1993 y 2008. Elaboración propia.}
  \label{fig:curvas}
\end{figure}

Dos lecturas inmediatas de la figura, que las secciones siguientes desarrollan.

\textbf{Primera: el cruce ocurrió en 2002, no en 2012.} Wikipedia pasó de
\num{10000} artículos en octubre de 2001 a \num{100000} en enero de 2003, de
modo que en algún momento de 2002 superó los \num{65000} de la Micropædia y los
\num{62000} de Encarta. La superioridad técnica en cobertura cambió de manos
\textbf{diez años antes} de que la impresión terminara y siete antes de que
Encarta cerrara. El retraso entre el cruce del FOM y la muerte comercial es la
medida de la inercia del canal de distribución.

\textbf{Segunda: la curva de Wikipedia ya saturó.} El tramo posterior a 2010 es
visiblemente plano en escala logarítmica: pasar de 6 a 7 millones de artículos
tomó cinco años y cuatro meses, frente a los diecinueve meses que tomó pasar de
1 a 2 millones. Una tecnología en la parte alta de su curva en S es
precisamente el objetivo natural del siguiente invasor.

\section{La tecnología establecida: la enciclopedia impresa}

\subsection*{01. El desempeño que la hizo dominante}

La decimoquinta edición, publicada en 1974 y reorganizada en 1985, costó
32 millones de dólares en trabajo editorial, la mayor inversión privada en
publicación de su época, y reunió a más de \num{4000} autores de más de cien
países. Su estructura final ofrecía unos \num{65000} artículos breves en la
Micropædia y unos setecientos tratados extensos en la Macropædia, con más de
40 millones de palabras en 32 volúmenes \parencite{britannica3,evans1997}.

Esa cifra define el FOM en el que la impresa era imbatible: \textbf{profundidad
por artículo}, con unas 615 palabras de media, y autoridad verificable, con
autores identificados y responsables. Todavía en 2005, cuando ya había perdido
la cobertura, seguía ganando en exactitud: en una muestra de 42 artículos
científicos revisados por especialistas, \emph{Nature} contó 123 errores en la
Britannica frente a 162 en Wikipedia \parencite{giles2005}.

\subsection*{02. El desempeño que la condenó}

En 1936 la Britannica adoptó la política de \textbf{revisión continua}: en lugar
de esperar una edición nueva cada veinticinco años, reimprimía cada año y
revisaba al menos el diez por ciento de los artículos, de modo que cada entrada
se reconsideraba unas dos veces por década \parencite{ebwiki}. Fue una mejora
real de su FOM de latencia, de unos veinticinco años a unos cinco, y también el techo de su curva:
por debajo de ese valor el papel no podía bajar.

La economía era todavía más rígida. Un juego se vendía entre \num{1500} y
\num{2200} dólares según la encuadernación, y su coste marginal se
descomponía en unos 250 dólares de producción más entre 500 y 600 dólares de
comisión para el vendedor \parencite{evans1997}. El activo real de la compañía no
era el contenido: era una fuerza de venta directa a domicilio. En 1990, el año
de su máximo histórico, las ventas de juegos multivolumen alcanzaron unos
650 millones de dólares, con \num{120000} juegos colocados en Estados Unidos.

\begin{figure}[ht]
  \centering
  \begin{tikzpicture}
    \fomgrid{4.55}
    \fombar{4}{7.188}{cImpresa}{\rCob}{\num{65000}}
    \fombar{3}{5.326}{cImpresa}{\rPro}{$\approx$ 615}
    \fombar{2}{10.303}{cImpresa!75}{\rLat}{$1{,}6\times10^{8}$ (5 años)}
    \fombar{1}{4.673}{cImpresa!75}{\rAcc}{$\approx$ 120 (supuesto)}
    \fombar{0}{5.431}{cImpresa!75}{\rCos{copia}}{750 a 850}
  \end{tikzpicture}
  \caption{Perfil de figuras de mérito de la enciclopedia impresa
  (\emph{Encyclopædia Britannica}, edición vigente entre 1985 y 2010) sobre un
  eje logarítmico común. \notafom{} Las barras oscuras marcan los FOM en los que
  la tecnología era superior y las claras aquellos en los que era vulnerable. La
  profundidad es el cociente entre más de 40 millones de palabras y
  \num{65000} artículos \parencite{evans1997}; el coste marginal suma producción
  y comisión de venta; la latencia de acceso es un supuesto explícito de orden
  de magnitud, no un dato medido. Elaboración propia.}
  \label{fig:fomimpresa}
\end{figure}

La Figura~\ref{fig:fomimpresa} muestra el perfil completo. Leída de arriba a
abajo describe con precisión el punto ciego del titular: era excelente en las dos
primeras filas y estaba entre cinco y ocho órdenes de magnitud por detrás de lo
posible en las tres últimas.

\section{Primer invasor: Microsoft Encarta}

\subsection*{01. Cuándo y con qué desempeño entró}

Microsoft cerró en 1989 un acuerdo de derechos no exclusivos con los editores de
la \emph{Funk \& Wagnalls Encyclopedia}, una obra de segundo nivel, y en 1993
lanzó Encarta \parencite{smith2017}. El contenido era claramente inferior:
7 millones de palabras frente a más de 40 millones de la Britannica
\parencite{evans1997}. En su FOM propio, el titular ganaba por un factor de casi
seis.

En todos los demás la relación se invertía. Encarta se lanzó a 395 dólares,
bajó pronto a 99 y con frecuencia se incluía sin coste adicional en la compra de
un ordenador nuevo. Su coste marginal era de 1,50 dólares por disco
\parencite{evans1997}, unas quinientas veces menos que el del juego impreso, lo
que hacía económicamente racional regalarla. El resultado fue inmediato: unas
\num{350000} copias a finales de 1993 y más de un millón en su segundo año
\parencite{encartawiki}.

\begin{figure}[ht]
  \centering
  \begin{tikzpicture}
    \fomgrid{4.55}
    \fombar{4}{7.169}{cEncarta!75}{\rCob}{\num{62000} \tiny(2008)}
    \fombar{3}{4.734}{cEncarta!75}{\rPro}{$\approx$ 140}
    \fombar{2}{9.659}{cEncarta}{\rLat}{$3{,}2\times10^{7}$ (1 año)}
    \fombar{1}{3.403}{cEncarta}{\rAcc}{$\approx$ 5}
    \fombar{0}{2.922}{cEncarta}{\rCos{copia}}{1,50}
    % marcadores de la posición del titular
    \draw[cImpresa, thick, densely dotted] (7.188,3.78) -- (7.188,4.22);
    \draw[cImpresa, thick, densely dotted] (5.326,2.78) -- (5.326,3.22);
    \draw[cImpresa, thick, densely dotted] (10.303,1.78) -- (10.303,2.22);
    \draw[cImpresa, thick, densely dotted] (4.673,0.78) -- (4.673,1.22);
    \draw[cImpresa, thick, densely dotted] (5.431,-0.22) -- (5.431,0.22);
  \end{tikzpicture}
  \caption{Perfil de figuras de mérito de Encarta. \notafom{} Las líneas
  punteadas azules marcan la posición de la enciclopedia impresa en cada FOM.
  Encarta pierde en las dos primeras filas y gana entre uno y tres órdenes de
  magnitud en las tres últimas. La profundidad es el cociente entre 7 millones
  de palabras
  \parencite{evans1997} y unos \num{50000} artículos de la edición estándar
  \parencite{encartawiki}; las
  dos cifras no proceden del mismo año, por lo que el valor es aproximado.
  Elaboración propia.}
  \label{fig:fomencarta}
\end{figure}

\subsection*{02. La reacción del titular y por qué no bastó}

Britannica no se quedó quieta, y esto merece subrayarse porque contradice la
caricatura del titular pasivo. Publicó su propio CD-ROM en 1994 a 995 dólares y
lo bajó a 200 en 1995, a 150 en 1997 y a 85 en 1998; ese mismo 1994 lanzó
Britannica Online, la primera enciclopedia de internet, con suscripciones de 150
y luego 85 dólares anuales; y en 1999 abrió britannica.com con el texto completo
gratuito y financiado con publicidad \parencite{smith2017}.

Reaccionó en producto y en precio con notable rapidez. Lo que no podía hacer era
sostener una comisión de 500 a 600 dólares por venta con un producto de 99
dólares. El obstáculo no era técnico ni cognitivo, era estructural: el canal que
había construido su dominio era incompatible con el nuevo punto de precio, que
es exactamente el mecanismo descrito por \textcite{christensen1997}. Las ventas
de enciclopedias impresas en Estados Unidos cayeron más del ochenta por ciento
desde 1990; en mayo de 1995 la Benton Foundation puso la compañía en venta y en
1996 el financiero Jacob Safra la compró por menos de la mitad de su valor en
libros \parencite{evans1997}.

\begin{quote}
\textbf{Consecuencia para el caso.} El golpe mortal a la enciclopedia impresa lo
dio el CD-ROM entre 1993 y 1996, cuando Wikipedia todavía no existía. Wikipedia
llegó a un mercado que ya estaba destruido y su víctima real fue Encarta.
\end{quote}

\section{Segundo invasor: Wikipedia}

\subsection*{01. El experimento de control: Nupedia}

Antes de Wikipedia, los mismos fundadores intentaron lo contrario. Nupedia
salió en línea el 9 de marzo de 2000 con un proceso editorial de siete pasos,
revisión por pares y expertos acreditados, es decir replicando el FOM del
titular. Produjo 24 artículos aprobados en tres años y medio, más unos 150 en
curso, y su servidor cayó definitivamente el 26 de septiembre de 2003
\parencite{nupediawiki}.

El contraste es el mejor argumento del caso: en ese mismo periodo, Wikipedia,
lanzada el 15 de enero de 2001, pasó de cero a más de \num{100000} artículos. La
diferencia no estaba en el contenido, en la financiación ni en las personas, que
eran las mismas. Estaba en el \textbf{coste de coordinación}: el wiki eliminó la
necesidad de asignar, aprobar y sincronizar el trabajo antes de publicarlo, lo
que \textcite{benkler2002} describe como producción entre pares, una forma de
organizar el esfuerzo que no es ni la empresa ni el mercado.

\subsection*{02. El desempeño actual}

\begin{figure}[ht]
  \centering
  \begin{tikzpicture}
    \fomgrid{4.55}
    \fombar{4}{9.071}{cWiki}{\rCob}{\num{7224454}}
    \fombar{3}{5.389}{cWiki}{\rPro}{$\approx$ 720}
    \fombar{2}{4.396}{cWiki}{\rLat}{$\approx$ 60 (minutos)}
    \fombar{1}{3.403}{cWiki}{\rAcc}{$\approx$ 5}
    \fombar{0}{0.049}{cWiki}{\rCos{consulta}}{$\approx$ 0,001}
    \draw[cImpresa, thick, densely dotted] (7.188,3.78) -- (7.188,4.22);
    \draw[cImpresa, thick, densely dotted] (5.326,2.78) -- (5.326,3.22);
    \draw[cImpresa, thick, densely dotted] (10.303,1.78) -- (10.303,2.22);
    \draw[cImpresa, thick, densely dotted] (4.673,0.78) -- (4.673,1.22);
    \draw[cImpresa, thick, densely dotted] (5.431,-0.22) -- (5.431,0.22);
    \draw[cEncarta, thick, densely dotted] (7.169,3.78) -- (7.169,4.22);
    \draw[cEncarta, thick, densely dotted] (4.734,2.78) -- (4.734,3.22);
    \draw[cEncarta, thick, densely dotted] (9.659,1.78) -- (9.659,2.22);
    \draw[cEncarta, thick, densely dotted] (3.403,0.78) -- (3.403,1.22);
    \draw[cEncarta, thick, densely dotted] (2.922,-0.22) -- (2.922,0.22);
    \node[font=\tiny, cImpresa, above=0.6mm] at (10.303,2.22) {impresa};
    \node[font=\tiny, cEncarta, below=0.6mm] at (9.659,1.78) {Encarta};
  \end{tikzpicture}
  \caption{Perfil de figuras de mérito de Wikipedia en agosto de 2026.
  \notafom{} Las líneas punteadas azules marcan la posición de la enciclopedia
  impresa y las naranjas la de Encarta. Es la primera
  tecnología de la cadena que gana en las cinco filas a la vez. La profundidad
  es el cociente entre más de \num{5000} millones de palabras
  \parencite{pew2026} y \num{7224454} artículos \parencite{wikistats2026}; el
  coste marginal es el gasto anual de la Fundación Wikimedia, 190,9 millones de
  dólares \parencite{wmf2025}, dividido entre las visitas anuales estimadas a
  partir de un promedio de 508 millones diarias \parencite{pew2026}. Elaboración
  propia.}
  \label{fig:fomwiki}
\end{figure}

Wikipedia es el único eslabón de la cadena que llegó a ganar en todos los FOM
del cuadro, incluida la profundidad: sus 720 palabras por artículo superan las
615 de la Britannica, resultado que en 2001 nadie habría anticipado. La brecha
de exactitud, el último refugio del titular, se había estrechado ya en 2005 a
una diferencia de 3,9 frente a 2,9 errores por artículo \parencite{giles2005}.

La escala actual es de otro orden: más de 66 millones de artículos en 342
idiomas, más de \num{600000} usuarios activos, unos \num{500} artículos nuevos al
día solo en inglés y una media de 5,7 ediciones por segundo en 2024
\parencite{pew2026,wikistats2026}. Y todo ello con un gasto anual de 190,9
millones de dólares \parencite{wmf2025}, menos de la tercera parte de lo que la
Britannica facturaba en 1990 en un solo país.

Microsoft cerró Encarta el 31 de octubre de 2009 con una explicación que es un
manual de gestión tecnológica en dos líneas: \enquote{la categoría de las
enciclopedias tradicionales y del material de referencia ha cambiado} y
\enquote{hoy la gente busca y consume información de maneras
considerablemente distintas a las del pasado} \parencite{techcrunch2009}.

\section{Tercer invasor: los modelos de lenguaje}

\subsection*{01. Cuándo entró y qué FOM redefine}

ChatGPT se publicó el 30 de noviembre de 2022, alcanzó 100 millones de usuarios
en dos meses y 800 millones de usuarios semanales en octubre de 2025
\parencite{altman2025}. La novedad, en términos de FOM, no es la cobertura ni el
coste: es que \textbf{deja de entregar el artículo y entrega la respuesta}. El
lector ya no tiene que localizar la entrada correcta, leerla y extraer de ella lo
que preguntaba. Cuando eso ocurre, el número de artículos deja de ser la unidad
relevante y la comparación en las mismas unidades se vuelve imposible, que es la
segunda propiedad enunciada en el marco.

El caso Grokipedia permite, por una vez, medirlo. Lanzada el 27 de octubre de
2025 con unas \num{885279} entradas generadas por un modelo,
\textcite{yasseri2025} la comparó con Wikipedia y encontró artículos
sistemáticamente más largos, \num{14200} palabras de media frente a \num{9400},
con una similitud semántica de 0,825 y, a la vez, 21,4 referencias por cada
\num{1000} palabras frente a 90,2, y 16,8 enlaces frente a 428. Su conclusión es
que la obra \enquote{reempaqueta en gran medida contenido humano ya curado a
través de una lente de inteligencia artificial que privilegia la fluidez y la
narrativa sobre la atribución}.

\begin{figure}[ht]
  \centering
  \begin{tikzpicture}
    \fomgrid{4.55}
    \fill[cLLM] (0,3.83) rectangle (8.231,4.17);
    \draw[-{Stealth[length=1.8mm]}, thick, cLLM] (8.30,4.0) -- (9.55,4.0);
    \node[left, text width=3.5cm, align=right, font=\scriptsize] at (-0.12,4)
      {\rCob};
    \node[right, font=\scriptsize] at (9.62,4) {\num{885279}, no acotada};
    \fombar{3}{6.580}{cLLM}{\rPro}{\num{14200}}
    \node[left, text width=3.5cm, align=right, font=\scriptsize] at (-0.12,2)
      {\rLat};
    \fill[pattern=north east lines, pattern color=cLLM]
      (2.76,1.83) rectangle (9.38,2.17);
    \draw[cLLM] (2.76,1.83) rectangle (9.38,2.17);
    \node[right, font=\scriptsize] at (9.48,2) {\tiny según la fuente};
    \fombar{1}{3.403}{cLLM}{\rAcc}{$\approx$ 5, ya sintetizada}
    \node[left, text width=3.5cm, align=right, font=\scriptsize] at (-0.12,0)
      {\rCos{consulta}};
    \node[right, font=\scriptsize, gray!45!black] at (0.1,0)
      {no divulgado por los proveedores};
    \draw[cWiki, thick, densely dotted] (9.071,3.78) -- (9.071,4.22);
    \draw[cWiki, thick, densely dotted] (5.389,2.78) -- (5.389,3.22);
    \draw[cWiki, thick, densely dotted] (3.403,0.78) -- (3.403,1.22);
  \end{tikzpicture}

  \vspace{0.4em}
  \begin{tikzpicture}
    \path (-3.62,-0.95) (12.20,3.55);
    \node[font=\scriptsize\itshape, gray!35!black, right] at (-3.62,3.35)
      {El FOM que empeora: fiabilidad de las respuestas};
    \foreach \xp/\lab in {0/0, 1.875/25, 3.75/50, 5.625/75, 7.5/100} {
      \draw[gray!25] (\xp,-0.42) -- (\xp,2.92);
      \node[below, font=\tiny, gray!55!black] at (\xp,-0.42) {\lab\,\%};
    }
    \foreach \y/\xv/\v/\lab in {%
      2.5/6.075/81/{Con algún problema},
      1.667/3.375/45/{Con al menos un\\problema significativo},
      0.833/2.325/31/{Con problemas graves\\de atribución de fuentes},
      0/1.500/20/{Con errores importantes\\de exactitud}} {
      \node[left, text width=3.45cm, align=right, font=\scriptsize]
        at (-0.15,\y) {\lab};
      \fill[cLLM!80] (0,\y-0.16) rectangle (\xv,\y+0.16);
      \node[right, font=\scriptsize] at (\xv+0.12,\y) {\v\,\%};
    }
  \end{tikzpicture}
  \caption{Perfil de figuras de mérito de los modelos de lenguaje como
  sustituto de la enciclopedia. \notafom{} Las líneas punteadas verdes marcan la
  posición de Wikipedia, que en la latencia de actualización queda dentro del
  intervalo rayado. Panel
  superior: los valores medibles proceden de Grokipedia \parencite{yasseri2025};
  la latencia de actualización se hereda de la fuente consultada, desde segundos
  cuando el modelo busca en la web hasta la fecha de corte de su entrenamiento
  cuando no lo hace. Panel inferior: el FOM que empeora. Corresponde al estudio
  de la \textcite{ebu2025} sobre \num{3000} respuestas en catorce idiomas.
  Elaboración propia.}
  \label{fig:fomllm}
\end{figure}

\subsection*{02. La paradoja del invasor parásito}

Este tercer invasor tiene una particularidad que no se dio en los dos anteriores:
\textbf{depende del titular al que desplaza}. Wikipedia aportó \num{3000}
millones de tokens al entrenamiento de GPT-3, apenas el 0,6 por ciento del corpus
en bruto, pero se le asignó el 3 por ciento de la mezcla de entrenamiento, de
modo que el modelo recorrió ese corpus 3,4 veces frente a las 0,44 veces que
recorrió Common Crawl: fue, con diferencia, la fuente más sobremuestreada del
conjunto \parencite{brown2020}. Y sigue siéndolo en la fase de uso: en un
análisis de unos 680 millones de citas de sistemas de IA entre agosto de 2024 y
junio de 2025, Wikipedia representó el 47,9 por ciento de los diez dominios más
citados por ChatGPT \parencite{profound2025}.

Al mismo tiempo, la sustitución ya se mide en el tráfico. Entre marzo y agosto de
2025 las visitas humanas a Wikipedia cayeron alrededor del 8 por ciento respecto
al mismo periodo de 2024, y la Fundación Wikimedia atribuye la caída a los
resúmenes generados por IA en los buscadores y al vídeo social
\parencite{miller2025}. El mecanismo está cuantificado: con un resumen de IA en
la página de resultados, los usuarios hacen clic en algún enlace en el 8 por
ciento de las visitas, frente al 15 por ciento cuando no lo hay, y solo en el 1
por ciento pulsan un enlace del propio resumen \parencite{pew2025}.

La respuesta del titular reproduce el patrón de 1994. Wikipedia adoptó el 4 de
agosto de 2025 un criterio de borrado rápido, el G15, para artículos claramente
generados por un modelo sin revisión humana \parencite{wikig15}, y la Fundación
monetiza al invasor a través de Wikimedia Enterprise, que facturó 8,3 millones de
dólares en el ejercicio 2024--2025, un 148 por ciento más que el anterior
\parencite{wmf2025}. Es la misma jugada que el CD-ROM de 1994: defender el activo
y cobrar por el acceso, sin poder cambiar la unidad de medida que el mercado
empieza a usar.

\section{¿Cuándo y por qué ocurrió?}

El Cuadro~\ref{tab:crono} responde al \enquote{cuándo} y el
Cuadro~\ref{tab:matriz} al \enquote{por qué}, comparando las cuatro
generaciones en el mismo conjunto de FOM.

\begin{table}[ht]
  \centering
  \small
  \caption{Cronología verificada del proceso de sustitución.}
  \label{tab:crono}
  \begin{tabularx}{\textwidth}{@{}l >{\raggedright\arraybackslash}X@{}}
    \toprule
    \textbf{Fecha} & \textbf{Hecho} \\
    \midrule
    1768 & Primera edición de la \emph{Encyclopædia Britannica}, en Edimburgo \\
    1936 & Política de revisión continua: la latencia baja de unos 25 años a
      unos 5 \\
    1974 y 1985 & Decimoquinta edición y su reorganización: \num{65000}
      artículos, 32 millones de dólares de inversión editorial \\
    1990 & Máximo histórico: unos 650 millones de dólares de ventas y
      \num{120000} juegos en Estados Unidos \\
    1993 & Lanzamiento de Encarta a 395 dólares, luego 99 y con frecuencia
      incluida con el ordenador \\
    1994 & Britannica responde con CD-ROM a 995 dólares y con Britannica
      Online, la primera enciclopedia de internet \\
    1995 y 1996 & La compañía se pone en venta y se vende por menos de la mitad
      de su valor en libros \\
    1999 & britannica.com se abre gratis con publicidad \\
    2000 & Nupedia: 24 artículos en tres años y medio con revisión de siete
      pasos \\
    15/01/2001 & Lanzamiento de Wikipedia \\
    2002 & \textbf{Cruce del FOM de cobertura}: Wikipedia supera los
      \num{65000} artículos de la Micropædia \\
    2005 & \emph{Nature}: 162 errores en Wikipedia frente a 123 en Britannica
      en 42 artículos científicos \\
    31/10/2009 & Cierre de Encarta \\
    13/03/2012 & Fin de la edición impresa de la Britannica, tras 244 años \\
    30/11/2022 & Lanzamiento de ChatGPT \\
    28/05/2025 & Wikipedia en inglés alcanza 7 millones de artículos \\
    2025 & Caída del 8\,\% en visitas humanas a Wikipedia; criterio G15 contra
      contenido generado por IA; Grokipedia con \num{885279} entradas \\
    \bottomrule
  \end{tabularx}
\end{table}

\begin{table}[ht]
  \centering
  \small
  \caption{Matriz de figuras de mérito de las cuatro generaciones.}
  \label{tab:matriz}
  \begin{tabularx}{\textwidth}{@{}>{\raggedright\arraybackslash}p{0.20\textwidth}
    >{\raggedright\arraybackslash}X >{\raggedright\arraybackslash}X
    >{\raggedright\arraybackslash}X >{\raggedright\arraybackslash}X@{}}
    \toprule
    \textbf{FOM} & \textbf{Impresa} & \textbf{Encarta} & \textbf{Wikipedia}
      & \textbf{Modelos de lenguaje} \\
    \midrule
    Cobertura (artículos)
      & \num{65000} & \num{62000} & \num{7224454} & no acotada \\
    Profundidad (palabras por artículo)
      & $\approx$ 615 & $\approx$ 140 & $\approx$ 720 & \num{14200} \\
    Latencia de actualización
      & 5 años & 1 año & minutos & según la fuente \\
    Latencia de acceso
      & minutos & segundos & segundos & segundos, ya sintetizada \\
    Coste para el lector
      & \num{1500} a \num{2200} USD & 99 USD o incluida & 0 & 0 o suscripción \\
    Coste marginal
      & 750 a 850 USD & 1,50 USD & $\approx$ 0,001 USD & no divulgado \\
    Idiomas
      & 1 & varios & 342 & prácticamente cualquiera \\
    Verificabilidad
      & autor identificado & editorial & historial y referencias públicas
      & 21,4 referencias por \num{1000} palabras \\
    Modelo de producción
      & \num{4000} autores contratados & licencia de contenido de terceros
      & producción entre pares & síntesis a partir de las tres anteriores \\
    \bottomrule
  \end{tabularx}
\end{table}

Cinco mecanismos explican el \enquote{por qué}, y ninguno de ellos es la calidad
del contenido.

\textbf{01. El coste marginal cae hasta desaparecer.} De 750 a 850 dólares por
copia impresa a 1,50 en CD y a una milésima de dólar por consulta en la web. Un
coste marginal casi nulo no solo abarata: cambia lo que es racional hacer con el
producto, y regalar la enciclopedia para vender ordenadores pasa a ser una buena
decisión.

\textbf{02. El soporte fija el techo de la latencia.} El papel no podía bajar de
unos años, el CD no podía bajar de una edición anual, la web bajó a minutos y el
modelo de lenguaje responde en segundos. Cada generación heredó de su soporte un
límite que ninguna mejora incremental podía superar, que es la definición de
techo de curva en S \parencite{foster1986}.

\textbf{03. El cuello de botella era coordinar, no escribir.} Nupedia y
Wikipedia tuvieron el mismo objetivo, el mismo fundador y diez meses de
diferencia entre sus lanzamientos. La diferencia de resultado, 24 artículos frente a más de \num{100000},
mide el coste de la coordinación jerárquica frente a la producción entre pares
\parencite{benkler2002}.

\textbf{04. El modelo de negocio ancla al titular.} Britannica podía fabricar un
CD-ROM, y lo hizo antes de que Encarta fuera dominante, pero no podía pagar una
comisión de 500 dólares con un producto de 99. El titular no perdió por no ver
la tecnología: perdió porque adoptarla implicaba destruir su propio canal
\parencite{christensen1997}.

\textbf{05. El invasor cambia la unidad de medida.} El paso decisivo no es
mejorar el FOM del titular, sino sustituirlo. Encarta cambió \enquote{palabras
por obra} por \enquote{coste de acceso}; Wikipedia cambió \enquote{autoridad del
autor} por \enquote{cobertura y actualidad}; los modelos de lenguaje cambian
\enquote{artículos disponibles} por \enquote{respuestas a mi pregunta}. Cuando
el titular ya no puede puntuar en el nuevo eje, la discusión sobre calidad deja
de ser relevante para el mercado.

\section{Precisiones y limitaciones}

\begin{enumerate}
  \item \textbf{Los FOM no son estrictamente conmensurables entre
    generaciones.} Un artículo de la Micropædia, uno de Wikipedia y una entrada
    generada por un modelo no son la misma unidad de producto, aunque compartan
    el nombre. Las comparaciones de las Figuras \ref{fig:fomimpresa} a
    \ref{fig:fomllm} son válidas como órdenes de magnitud y no como medidas
    exactas.
  \item \textbf{Dos valores del cuadro son supuestos declarados.} La latencia de
    acceso de la enciclopedia impresa, unos 120 segundos, es una estimación de
    orden de magnitud, y la latencia de actualización de los modelos de lenguaje
    es un intervalo, no un valor, porque depende de si el sistema consulta la web
    en el momento de responder.
  \item \textbf{Dos tramos de la Figura~\ref{fig:curvas} son esquemáticos.} No
    existe una serie anual pública del recuento de artículos de Encarta entre
    1993 y 2008, ni una serie homogénea de la Britannica anterior a 1889.
  \item \textbf{El efecto vela es un fenómeno discutido.} Es tentador presentar
    la decimoquinta edición de 1974, la mayor inversión editorial de su historia,
    o el CD-ROM de 1994, como una aceleración defensiva del titular ante la
    amenaza. Conviene ser prudente: \textcite{howells2002} reexaminó los casos
    clásicos, la vela frente al vapor y el proceso Leblanc frente al Solvay, y
    concluyó que el efecto no existió en ninguno de los dos, y
    \textcite{mendonca2013} matiza que lo que suele observarse es una mejora
    continua ya en marcha, no una reacción. Aplicado a este caso: la revisión
    continua de 1936 y la decimoquinta edición son anteriores a la amenaza
    digital y no pueden leerse como respuesta a ella.
  \item \textbf{La cuarta sustitución está en curso, no consumada.} Wikipedia no
    ha muerto: sigue creciendo unos \num{500} artículos diarios en inglés y es la
    fuente más citada por los propios modelos que la desplazan. Lo que cae es su
    tráfico humano, no su corpus. Los datos disponibles describen el comienzo de
    una invasión, no su desenlace.
  \item \textbf{La tecnología murió, la empresa no.} Britannica abandonó la
    imprenta y se convirtió en un proveedor de software educativo con IA, con
    ingresos que apuntaban a los 100 millones de dólares y una valoración buscada
    de unos mil millones en una salida a bolsa \parencite{bloomberg2024,
    delamerced2024}. Sobrevivir a la muerte de la propia tecnología exige
    cambiar de curva, y no todos los titulares de esta historia lo consiguieron:
    Encarta no cambió de curva y desapareció.
\end{enumerate}

\section*{Conclusiones}

La enciclopedia impresa fue reemplazada, y el proceso admite una respuesta
precisa a las dos preguntas del enunciado.

\textbf{Cuándo.} La superioridad técnica cambió de manos en dos momentos
distintos, ninguno de los cuales coincide con el hecho que suele citarse. En
coste y accesibilidad, entre 1993 y 1996, cuando el CD-ROM redujo el coste
marginal de una copia de unos 800 dólares a 1,50 y hundió las ventas de la
impresa más del ochenta por ciento. En cobertura, en 2002, cuando Wikipedia
superó los \num{65000} artículos de la Micropædia. El fin de la impresión, en
marzo de 2012, es el último acto administrativo de una sustitución que se había
decidido diez años antes.

\textbf{Por qué.} Porque el titular era excelente en las dos figuras de mérito
que él mismo consideraba decisivas, profundidad y autoridad, y estaba entre cinco
y ocho órdenes de magnitud por detrás de lo posible en las tres que acabaron
decidiendo el mercado: latencia de actualización, latencia de acceso y coste
marginal. Y porque esas tres no podían mejorarse sin destruir el canal de venta
directa que constituía su verdadero activo.

La lección que me llevo como CTO no es que convenga vigilar a los invasores, que
es obvio, sino algo más incómodo: \textbf{el FOM en el que uno es mejor suele ser
el que el mercado está a punto de dejar de usar}. La Britannica ganaba en
palabras por artículo y perdió en dólares por consulta. Wikipedia ganó en
artículos accesibles y hoy compite contra sistemas a los que el recuento de
artículos no se les puede aplicar. En una cadena de sustituciones, quien mide su
posición solo con el eje heredado ya está midiendo la carrera equivocada.

\printbibliography[title={Referencias bibliográficas}]

\end{document}
